% backmatter/glossary.tex
%
% Глоссарий аббревиатур — двухколоночный longtable, разделён на латинский
% и~кириллический блоки. Записи добавляются в~алфавитном порядке внутри
% соответствующего блока; ссылок на~глоссарий из~корпуса нет (приложение,
% а~не индекс терминов).

\addchap{Глоссарий}\label{ch:glossary}

\begingroup
\footnotesize
\begin{longtable}{@{}p{4.5em}p{\dimexpr\linewidth-4.5em-2\tabcolsep\relax}@{}}
  \textbf{A2A} & Account-to-Account, перевод со счёта на счёт без участия карточной сети; общий случай account-based-модели платежа в~противоположность card-based-операциям карточных схем \\
  \textbf{AAC} & Application Authentication Cryptogram, криптограмма отказа в~EMV \\
  \textbf{ABA} & American Bankers Association, формат Track~2 магнитной полосы \\
  \textbf{ABU} & Automatic Billing Updater, сервис Mastercard автоматического обновления реквизитов карт у~мерчантов; аналог Visa Account Updater (см.~VAU) \\
  \textbf{AC} & Application Cryptogram, криптограмма EMV-операции (одна из ARQC, TC, AAC) \\
  \textbf{ACH} & Automated Clearing House, расчётная сеть пакетных межбанковских переводов в~США \\
  \textbf{ACS} & Access Control Server, компонент эмитента в~3-D~Secure, проводящий аутентификацию держателя \\
  \textbf{AES} & Advanced Encryption Standard, симметричный блочный шифр (FIPS 197) \\
  \textbf{AFL} & Application File Locator, описывает расположение записей EMV-приложения на чипе \\
  \textbf{AFT} & Account Funding Transaction, операция загрузки средств на счёт через карточную сеть (Visa Direct/Money Movement) \\
  \textbf{AID} & Application Identifier, идентификатор платёжного приложения на чипе карты (EMV) \\
  \textbf{AIP} & Application Interchange Profile, битовое поле, описывающее возможности EMV-приложения \\
  \textbf{AIS} & Account Information Security, программа Visa, предшественница PCI~DSS \\
  \textbf{AISP} & Account Information Service Provider, поставщик услуг информирования о~счёте (PSD2) \\
  \textbf{AML} & Anti-Money Laundering, противодействие отмыванию доходов \\
  \textbf{ANSI} & American National Standards Institute, американский институт национальных стандартов \\
  \textbf{AOC} & Attestation of Compliance, аттестат соответствия PCI~DSS \\
  \textbf{AOV} & Average Order Value, средний чек заказа \\
  \textbf{APDU} & Application Protocol Data Unit, единица обмена данными между терминалом и~чипом (ISO/IEC 7816-4) \\
  \textbf{API} & Application Programming Interface, программный интерфейс взаимодействия систем \\
  \textbf{APID} & Application Program Identifier, объект EMV (тег \texttt{0x9F5A}) в~бесконтактных профилях Visa; кодирует регион программы и~служит ключом выбора набора лимитов в~DRL \\
  \textbf{Apple ECP} & Enhanced Contactless Polling, проприетарное расширение Apple над ISO/IEC~14443-A для считывателей NFC (используется в~Apple Pay, Apple VAS, Apple Wallet Access) \\
  \textbf{ARCA} & Armenian Card, армянская национальная платёжная система \\
  \textbf{ARN} & Acquirer Reference Number, эквайринговый идентификатор операции, используемый в~клиринге и~спорах \\
  \textbf{ARPC} & Authorisation Response Cryptogram, ответная криптограмма эмитента (EMV online auth) \\
  \textbf{ARQC} & Authorisation Request Cryptogram, запросная криптограмма карты (EMV online auth) \\
  \textbf{ASK} & Amplitude Shift Keying, амплитудно-манипулированная модуляция в~ISO/IEC~14443 Type~A \\
  \textbf{ASPSP} & Account Servicing Payment Service Provider, банк, обслуживающий счёт клиента в~открытом банкинге PSD2 \\
  \textbf{ASV} & Approved Scanning Vendor, аккредитованный PCI~SSC внешний сканировщик уязвимостей \\
  \textbf{ATC} & Application Transaction Counter, счётчик операций EMV-приложения \\
  \textbf{ATM} & Automated Teller Machine, банкомат \\
  \textbf{ATO} & Account Takeover, захват учётной записи \\
  \textbf{AUC} & Area Under Curve, площадь под ROC- или PR-кривой; метрика классификатора \\
  \textbf{AVS} & Address Verification Service, проверка адреса держателя у~эмитента в~операциях CNP \\
  \textbf{BCB} & Banco Central do Brasil, Центральный банк Бразилии (регулятор PIX) \\
  \textbf{BCD} & Binary Coded Decimal, двоично-десятичное кодирование (ISO~8583, EMV) \\
  \textbf{BDK} & Base Derivation Key, базовый ключ деривации в~DUKPT \\
  \textbf{BIN} & Bank Identification Number, первые цифры PAN, идентифицирующие эмитента (см.~также IIN) \\
  \textbf{BIS} & Bank for International Settlements, Банк международных расчётов \\
  \textbf{BLE} & Bluetooth Low Energy, маломощный Bluetooth для платежей и~смежных задач \\
  \textbf{BNPL} & Buy Now Pay Later, рассрочка/отложенная оплата как платёжный метод \\
  \textbf{BRAM} & Business Risk Assessment and Mitigation, программа Mastercard по оценке рисков мерчанта \\
  \textbf{C2B} & Customer-to-Business, перевод от физлица в~пользу бизнеса \\
  \textbf{C2C} & Customer-to-Customer, перевод между физическими лицами \\
  \textbf{CA} & Certificate Authority, удостоверяющий центр; в~платежах~-- корень доверия сертификатов EMV карт и~PKI для удалённой инъекции ключей \\
  \textbf{CAVV} & Cardholder Authentication Verification Value, криптографический результат 3-D~Secure-аутентификации (Visa) \\
  \textbf{CBC} & Cipher Block Chaining, режим сцепления блоков для блочных шифров \\
  \textbf{CBDC} & Central Bank Digital Currency, цифровая валюта центрального банка (например, цифровой рубль) \\
  \textbf{CBPII} & Card-Based Payment Instrument Issuer, регистрируемая роль эмитента карточного инструмента (PSD2) \\
  \textbf{CBPR} & Cross-Border Payments and Reporting Plus, миграционная программа SWIFT на ISO~20022 \\
  \textbf{CDA} & Combined Data Authentication, метод офлайн-аутентификации данных карты с~динамической подписью операции (EMV) \\
  \textbf{CDCVM} & Consumer Device Cardholder Verification Method, верификация держателя средствами устройства (биометрия, PIN-код экрана) \\
  \textbf{CDE} & Cardholder Data Environment, периметр окружения данных держателя по PCI~DSS \\
  \textbf{CE} & Compelling Evidence, набор доказательств, позволяющий эквайеру отбить часть споров (Visa CE 3.0, MC compelling evidence) \\
  \textbf{CHD} & Cardholder Data, данные держателя карты в~терминах PCI~DSS \\
  \textbf{CIBA} & Client-Initiated Backchannel Authentication, протокол отделённой аутентификации (OpenID Foundation) \\
  \textbf{CIPS} & Cross-Border Interbank Payment System, китайская система международных межбанковских переводов в~юанях \\
  \textbf{CISP} & Cardholder Information Security Program, программа Visa, предшественница PCI~DSS \\
  \textbf{CIT} & Customer-Initiated Transaction, операция, инициированная держателем карты \\
  \textbf{CMAC} & Cipher-based Message Authentication Code, MAC на блочном шифре (NIST~SP~800-38B) \\
  \textbf{CMS} & Card Management System, система управления жизненным циклом карт у~эмитента \\
  \textbf{CNP} & Card-Not-Present, операция без физического предъявления карты \\
  \textbf{CPM} & Customer-Presented Mode, QR-сценарий с~предъявлением кода клиентом \\
  \textbf{CPS} & Custom Payment Service, программа квалификации операций Visa с~льготным interchange (CPS/Retail и~др.) \\
  \textbf{CUP} & China UnionPay, китайская карточная схема \\
  \textbf{CVC1} & Card Verification Code~1, статический код подлинности на магнитной полосе \\
  \textbf{CVC3} & Card Verification Code 3, динамический код Mastercard для contactless mag-stripe (PayPass M/Stripe) \\
  \textbf{CVM} & Cardholder Verification Method, метод верификации держателя (PIN, signature, CDCVM, no~CVM) \\
  \textbf{CVV/CVC} & Card Verification Value (Visa) / Card Verification Code (Mastercard), статический код подлинности карты на магнитной полосе \\
  \textbf{CVV2/CVC2} & Печатный код безопасности на оборотной стороне карты, используется в~операциях CNP \\
  \textbf{DCC} & Dynamic Currency Conversion, конверсия в~валюту держателя на стороне эквайера/терминала \\
  \textbf{DICT} & Diretório de Identificadores de Contas Transacionais, директория Banco Central do Brasil, отображающая ключи Pix (телефон, e-mail, CPF, CNPJ, EVP) в~реквизиты счёта и~участника \\
  \textbf{DDA} & Dynamic Data Authentication, метод офлайн-аутентификации с~динамической подписью данных карты (EMV) \\
  \textbf{DE} & Data Element, поле сообщения ISO 8583 \\
  \textbf{DLT} & Distributed Ledger Technology, технология распределённых реестров; база данных с~синхронизированными копиями у~нескольких участников \\
  \textbf{DMS} & Dual Message System, двухстадийная авторизация и~клиринг отдельным сообщением (см.~также SMS) \\
  \textbf{DPAN} & Device Primary Account Number, токенизированный номер, привязанный к~устройству (Apple Pay, Google Pay и~др.) \\
  \textbf{DRL} & Dynamic Reader Limit, механизм бесконтактных ридеров Visa и~AmEx, позволяющий держать несколько наборов лимитов на один AID в~зависимости от региональной принадлежности карты (см.~APID) \\
  \textbf{DS} & Directory Server, маршрутизатор сообщений 3-D~Secure от мерчанта к~ACS эмитента \\
  \textbf{DSRP} & Digital Secure Remote Payment, механизм Mastercard генерации криптограмм в~CNP \\
  \textbf{DUKPT} & Derived Unique Key Per Transaction, схема производных одноразовых ключей для PIN/данных в~POS \\
  \textbf{EAS} & Estimated Authorization Service, схема предварительных авторизаций с~уточнением суммы (Visa) \\
  \textbf{EBA} & European Banking Authority, Европейская банковская администрация \\
  \textbf{EBCDIC} & Extended Binary Coded Decimal Interchange Code, кодировка IBM, исторически используется в~ISO~8583 \\
  \textbf{ECC} & Elliptic Curve Cryptography, криптография на эллиптических кривых \\
  \textbf{ECDSA} & Elliptic Curve Digital Signature Algorithm, алгоритм цифровой подписи на эллиптических кривых \\
  \textbf{ECI} & Electronic Commerce Indicator, индикатор результата 3-D~Secure в~авторизационном сообщении \\
  \textbf{ECM} & Excessive Chargeback Merchant, статус мониторинговой программы Mastercard для ТСП с~повышенной долей chargeback \\
  \textbf{EEA} & European Economic Area, Европейская экономическая зона \\
  \textbf{EFA} & EMV Final Advice, индустриальное название второй POS-to-host-сессии (ISO~8583 MTI~0220/0230), в~которой терминал передаёт ПЦ результат обработки картой ARPC и~финальную криптограмму; в~спецификациях EMVCo тот же шаг обозначается как online completion / second GENERATE~AC \\
  \textbf{EFM} & Excessive Fraud Merchant, статус мониторинговой программы Mastercard по фроду \\
  \textbf{EIRF} & Electronic Interchange Reimbursement Fee, штрафная категория interchange Visa для операций без полного набора данных \\
  \textbf{ELCARD} & Кыргызская национальная платёжная система (см.~также ЭЛКАРТ) \\
  \textbf{EMV} & Europay/Mastercard/Visa, семейство спецификаций чиповых карт (EMVCo) \\
  \textbf{EP} & E-commerce Page, тип SAQ A-EP в~PCI~DSS для платёжных страниц с~прямым постом данных карты \\
  \textbf{FAPI} & Financial-grade API, профиль безопасности OpenID Foundation для финансовых API (см.~также ФАПИ) \\
  \textbf{FATF} & Financial Action Task Force, межправительственная организация по борьбе с~отмыванием денег и~финансированием терроризма \\
  \textbf{FBO} & Fulfillment by Operator, маркетплейсная схема с~логистикой на стороне оператора \\
  \textbf{FBS} & Fulfillment by Seller, маркетплейсная схема с~логистикой на стороне продавца \\
  \textbf{FCA} & Financial Conduct Authority, регулятор финансовых услуг Великобритании \\
  \textbf{FedNow} & Мгновенная розничная платёжная система Federal Reserve Banks, запущена 20 июля 2023~года; регулируется 12~CFR~210 Subpart~C и~Operating Circular~8, расчёт через master account участника в~Reserve Bank \\
  \textbf{FF1} & Format-Preserving Encryption Algorithm~1, режим сохраняющего формат шифрования (NIST~SP~800-38G) \\
  \textbf{FF3} & Format-Preserving Encryption Algorithm~3, режим FPE; в~актуальной редакции FF3-1 (NIST~SP~800-38G Rev.~1) \\
  \textbf{FIN} & SWIFT FIN, исторический сервис передачи MT-сообщений \\
  \textbf{FIPS} & Federal Information Processing Standards, серия стандартов NIST (FIPS~140-3, FIPS~197 и~др.) \\
  \textbf{FIS} & Fidelity National Information Services, крупный американский процессор \\
  \textbf{Floor limit} & терминальный порог офлайн-одобрения по сумме операции; задаётся эквайером в~рамках правил схемы (EMV terminal risk management) \\
  \textbf{FPAN} & Funding Primary Account Number, реальный номер карты, лежащий за токеном \\
  \textbf{FPE} & Format-Preserving Encryption, шифрование с~сохранением формата (например, длины и~алфавита PAN) \\
  \textbf{FWT} & Frame Waiting Time, максимальное время ожидания ответа карты в~ISO/IEC~14443 \\
  \textbf{GCMS} & Global Clearing Management System, клиринговая система Mastercard \\
  \textbf{GPO} & Get Processing Options, EMV-команда инициации транзакции \\
  \textbf{HCE} & Host Card Emulation, эмуляция карты в~ОС устройства без аппаратного Secure Element \\
  \textbf{HECM} & High Excessive Chargeback Merchant, повышенный уровень программы Mastercard по chargeback \\
  \textbf{HMAC} & Hash-based Message Authentication Code, MAC на хеш-функции (RFC~2104) \\
  \textbf{HPP} & Hosted Payment Page, размещённая платёжная страница на стороне PSP \\
  \textbf{HSM} & Hardware Security Module, аппаратный модуль для хранения ключей и~криптоопераций \\
  \textbf{HUMO} & Узбекская национальная платёжная система \\
  \textbf{IAC} & Issuer Action Codes, заданные эмитентом правила решения по EMV-операции \\
  \textbf{IANA} & Internet Assigned Numbers Authority, реестр идентификаторов Интернета (порты, MIME, JOSE-алгоритмы) \\
  \textbf{IATA} & International Air Transport Association, формат Track~1 магнитной полосы \\
  \textbf{ICC} & Integrated Circuit Card, чиповая карта (EMV) \\
  \textbf{IFR} & Interchange Fee Regulation, регламент ЕС~2015/751 о~межбанковских комиссиях \\
  \textbf{IFS} & Interface Specification, документ карточной схемы, описывающий формат и~семантику ISO~8583-сообщений конкретного интерфейса \\
  \textbf{IIN} & Issuer Identification Number, термин ISO, эквивалент BIN \\
  \textbf{IK} & Initial Key, начальный ключ устройства в~DUKPT \\
  \textbf{IKI} & Initial Key Identifier, идентификатор начального ключа DUKPT \\
  \textbf{IMK} & Issuer Master Key, мастер-ключ эмитента, из которого деривируются ключи EMV-приложения карты \\
  \textbf{IPC} & Interbank Processing Center, межбанковский процессинговый центр Кыргызстана \\
  \textbf{IPR} & Instant Payments Regulation, Regulation (EU) 2024/886 о~мгновенных кредитовых переводах в~евро; обязательность SCT Inst в~еврозоне с~9~октября 2025~года \\
  \textbf{IPM} & Integrated Product Messages, формат клиринговых сообщений Mastercard \\
  \textbf{IRF} & Interchange Reimbursement Fee, межбанковская комиссия (термин UnionPay) \\
  \textbf{ISO} & International Organization for Standardization, в~книге упоминается через стандарты 4217, 7810, 7813, 7816, 8583, 20022 \\
  \textbf{JARM} & JWT Secured Authorization Response Mode, подписанный ответ авторизационного эндпоинта (FAPI) \\
  \textbf{JCB} & Japan Credit Bureau, японская карточная схема \\
  \textbf{JSON} & JavaScript Object Notation, формат обмена данными в~REST/HTTP-интерфейсах \\
  \textbf{JWE} & JSON Web Encryption (RFC~7516) \\
  \textbf{JWK} & JSON Web Key (RFC~7517) \\
  \textbf{JWS} & JSON Web Signature (RFC~7515) \\
  \textbf{JWT} & JSON Web Token (RFC~7519) \\
  \textbf{KCV} & Key Check Value, контрольная сумма ключа для верификации совпадения копий ключа на разных устройствах \\
  \textbf{KEK} & Key Encryption Key, ключ шифрования ключей; используется для защищённой передачи и~хранения других ключей \\
  \textbf{KYB} & Know Your Business, идентификация юридических лиц при подключении \\
  \textbf{KYC} & Know Your Customer, идентификация и~верификация клиента-физлица \\
  \textbf{LLVAR} & Length-Length Variable, поле переменной длины с~двухсимвольным префиксом длины (ISO~8583) \\
  \textbf{LMK} & Local Master Key, локальный мастер-ключ HSM; защищает производные ключи внутри устройства \\
  \textbf{MAC} & Message Authentication Code, код аутентичности сообщения \\
  \textbf{MC ECP} & Excessive Chargeback Program, программа Mastercard по мониторингу возвратных платежей у~эквайеров и~ТСП \\
  \textbf{MCC} & Merchant Category Code, четырёхзначный код категории мерчанта (ISO 18245) \\
  \textbf{MDB} & Mir Debit Business, продуктовый код дебетовой карты Мир для бизнеса \\
  \textbf{MDES} & Mastercard Digital Enablement Service, токенизационный сервис Mastercard \\
  \textbf{MDP} & Mir Debit Prepaid, продуктовый код предоплаченной карты Мир \\
  \textbf{MDR} & Merchant Discount Rate, торговая уступка мерчанта; синоним MSC \\
  \textbf{MED} & Mecanismo Especial de Devolução, возврат по антифроду Pix; банк-отправитель инициирует через инфраструктуру схемы запрос на блокировку и~возврат при подозрении на фрод \\
  \textbf{MDT} & Mir Debit, продуктовый код стандартной дебетовой карты Мир \\
  \textbf{Me2Me Pull} & Сценарий СБП, инициируемый банком-получателем: клиент даёт согласие в~приложении банка-получателя и~\enquote{стягивает} средства со своего счёта в~банке-отправителе \\
  \textbf{Me2Me Push} & Сценарий СБП, инициируемый банком-отправителем: клиент отправляет средства из приложения банка-отправителя на свой счёт в~банке-получателе \\
  \textbf{MID} & Merchant ID, идентификатор ТСП у~эквайера \\
  \textbf{MIT} & Merchant-Initiated Transaction, операция, инициированная мерчантом без участия держателя в~момент списания \\
  \textbf{MPA} & MIR Payment Application, EMV-апплет карты системы Мир с~поддержкой контактного и~бесконтактного интерфейсов и~криптографии ГОСТ \\
  \textbf{MPM} & Merchant-Presented Mode, QR-сценарий с~предъявлением кода продавцом \\
  \textbf{MSC} & Merchant Service Charge, торговая уступка, комиссия мерчанта эквайеру \\
  \textbf{MT} & Message Type, формат текстовых сообщений SWIFT (MT103, MT202 и~др.) \\
  \textbf{MTI} & Message Type Indicator, четырёхзначный идентификатор типа сообщения ISO 8583 \\
  \textbf{NFC} & Near Field Communication, ближнепольная радиосвязь для бесконтактных платежей \\
  \textbf{NIST} & National Institute of Standards and Technology, американский институт стандартов и~технологий \\
  \textbf{NNSS} & National Net Settlement Service, расчётный сервис Visa в~США \\
  \textbf{NPCI} & National Payments Corporation of India, оператор UPI и~RuPay в~Индии \\
  \textbf{NRZ} & Non-Return-to-Zero, схема линейного кодирования (например, в~ISO/IEC~14443 Type~B) \\
  \textbf{NTID} & Network Transaction ID, идентификатор операции сети Visa, используемый в~цепочках MIT/CIT и~CE 3.0 \\
  \textbf{OCT} & Original Credit Transaction, обратный денежный перевод на карту через карточную сеть \\
  \textbf{OFAC} & Office of Foreign Assets Control, санкционное подразделение Минфина США \\
  \textbf{OIDF} & OpenID Foundation, организация-разработчик OIDC и~FAPI \\
  \textbf{OXXO} & Мексиканская сеть convenience-магазинов и~одноимённый способ оплаты по штрихкоду \\
  \textbf{P2PE} & Point-to-Point Encryption, шифрование данных карты от точки приёма до дешифрующего узла; PCI имеет одноимённый стандарт \\
  \textbf{PAN} & Primary Account Number, основной номер карты (ISO/IEC 7812) \\
  \textbf{PAR} & Pushed Authorization Requests, режим OAuth с~предварительной отправкой параметров (RFC~9126) \\
  \textbf{PCI} & Payment Card Industry, отрасль платёжных карт; в~книге чаще про PCI~SSC и~его стандарты \\
  \textbf{PIN} & Personal Identification Number, секретный код держателя для верификации \\
  \textbf{PISP} & Payment Initiation Service Provider, поставщик услуг инициации платежа (PSD2) \\
  \textbf{PIX} & Бразильская система мгновенных платежей под управлением BCB \\
  \textbf{PKCE} & Proof Key for Code Exchange, защита OAuth-кода от перехвата (RFC~7636) \\
  \textbf{PKI} & Public Key Infrastructure, инфраструктура открытых ключей: удостоверяющие центры, сертификаты и~правила доверия между ними \\
  \textbf{PoA} & Proof-of-Authority, алгоритм консенсуса с~заранее утверждённым списком валидаторов; типичный выбор для приватных DLT и~CBDC \\
  \textbf{POI} & Point of Interaction, обобщённый термин PCI для устройств приёма платежа \\
  \textbf{POS} & Point of Sale, точка продажи; в~узком смысле -- POS-терминал \\
  \textbf{PoS} & Proof-of-Stake, алгоритм консенсуса, в~котором право подписать блок получают участники, заблокировавшие залог криптовалюты; нарушители теряют залог (\textit{slashing}). Применяется в~Ethereum с~2022~г. \\
  \textbf{PoW} & Proof-of-Work, алгоритм консенсуса, в~котором валидаторы тратят вычислительную работу на~подбор хеша; применяется в~Bitcoin \\
  \textbf{PPSE} & Proximity Payment System Environment, бесконтактный аналог PSE (\texttt{2PAY.SYS.DDF01}) \\
  \textbf{PSD2} & Payment Services Directive 2, директива ЕС об~открытом банкинге и~SCA \\
  \textbf{PSE} & Payment System Environment, контактный каталог приложений на чипе (\texttt{1PAY.SYS.DDF01}) \\
  \textbf{PSP} & Payment Service Provider, провайдер платёжных услуг \\
  \textbf{PSU} & Payment Service User, клиент в~терминологии PSD2 \\
  \textbf{PTS} & PIN Transaction Security, стандарт PCI~SSC для аппаратных модулей и~POS-устройств \\
  \textbf{PWCB} & Purchase with Cashback, тип EMV-операции \enquote{покупка с~выдачей наличных} (Mastercard); кодируется тегом \texttt{0x9C} = \texttt{09} \\
  \textbf{QMAP} & Questionable Merchant Audit Program, программа Visa по мониторингу подозрительных мерчантов \\
  \textbf{QR} & Quick Response (code), двумерный штрихкод как контейнер платёжных реквизитов \\
  \textbf{QSA} & Qualified Security Assessor, аккредитованный PCI~SSC аудитор \\
  \textbf{RBA} & Risk-Based Authentication, риск-ориентированная аутентификация в~3-D~Secure \\
  \textbf{RID} & Registered Application Provider Identifier, первая часть AID, идентифицирующая платёжную сеть (EMV) \\
  \textbf{RKI} & Remote Key Injection, удалённая загрузка криптографических ключей в~POS-терминал по~защищённому каналу \\
  \textbf{ROC} & Report on Compliance, итоговый отчёт QSA по PCI~DSS \\
  \textbf{RRN} & Retrieval Reference Number, идентификатор, по которому участники находят операцию (DE 37 в~ISO 8583) \\
  \textbf{RRP} & Relay-Resistant Protocol, защита от relay-атак в~Mastercard M/Chip \\
  \textbf{RSA} & Rivest-Shamir-Adleman, асимметричный криптоалгоритм \\
  \textbf{RTGS} & Real-Time Gross Settlement, валовой расчёт в~реальном времени \\
  \textbf{RTO} & Recovery Time Objective, целевое время восстановления сервиса после инцидента \\
  \textbf{RTS} & Regulatory Technical Standards, технические регламентирующие стандарты EBA по PSD2 \\
  \textbf{SAD} & Sensitive Authentication Data, данные аутентификации, не подлежащие хранению по PCI~DSS \\
  \textbf{SAM} & Switching and Account Management, коммутирующая инфраструктура Mastercard \\
  \textbf{SAQ} & Self-Assessment Questionnaire, опросник самооценки PCI~DSS \\
  \textbf{SCA} & Strong Customer Authentication, строгая аутентификация клиента (PSD2) \\
  \textbf{SCF} & Stored Credential Framework, правила Visa и~Mastercard для повторных списаний по~сохранённым реквизитам \\
  \textbf{SCT Inst} & SEPA Instant Credit Transfer, мгновенный кредитовый перевод в~евро по rulebook EPC; целевое end-to-end время 10~секунд, 24/7/365 \\
  \textbf{SDA} & Static Data Authentication, статическая офлайн-аутентификация данных карты (EMV) \\
  \textbf{SDK} & Software Development Kit, набор библиотек для интеграции в~приложение \\
  \textbf{SDN} & Specially Designated Nationals (and Blocked Persons), санкционный список OFAC \\
  \textbf{SDP} & Site Data Protection, программа Mastercard, предшественница PCI~DSS \\
  \textbf{SE} & Secure Element, аппаратный защищённый элемент устройства \\
  \textbf{SEPA} & Single Euro Payments Area, единое пространство платежей в~евро \\
  \textbf{SLA} & Service Level Agreement, соглашение об~уровне обслуживания \\
  \textbf{SMC} & Secure Messaging Confidentiality, ключ EMV для шифрования сценариев эмитента \\
  \textbf{SMI} & Secure Messaging Integrity, ключ EMV для целостности сценариев эмитента \\
  \textbf{SNA} & Systems Network Architecture, проприетарный сетевой стек IBM \\
  \textbf{SPI} & Sistema de Pagamentos Instantâneos, бразильская инфраструктура мгновенных расчётов под PIX \\
  \textbf{SRED} & Secure Reading and Exchange of Data, модуль PCI~PTS для шифрования данных карты внутри терминала \\
  \textbf{STAN} & System Trace Audit Number, локальный счётчик операций терминала/эквайера (DE 11 в~ISO 8583) \\
  \textbf{STIP} & Stand-In Processing, авторизация на стороне сети, когда эмитент недоступен \\
  \textbf{SWIFT} & Society for Worldwide Interbank Financial Telecommunication, межбанковская сеть передачи финансовых сообщений \\
  \textbf{TAC} & Terminal Action Codes, заданные сетью правила решения по EMV-операции на терминале \\
  \textbf{TADG} & Visa Transaction Acceptance Device Guide, документация интерфейса Visa для эквайера \\
  \textbf{TAVV} & Token Authentication Verification Value, криптограмма токена в~CNP-операциях (Visa Token Service) \\
  \textbf{TC} & Transaction Certificate, криптограмма операции EMV, одобренной офлайн \\
  \textbf{TC33} & Формат клирингового файла Visa BASE II \\
  \textbf{TDEA} & Triple Data Encryption Algorithm, он же Triple~DES (3DES); устаревший симметричный шифр \\
  \textbf{TDOL} & Transaction Data Object List, список тегов EMV для криптограммы~TC, задаваемый картой по умолчанию \\
  \textbf{TID} & Terminal ID, идентификатор терминала у~эквайера \\
  \textbf{TIPS} & TARGET Instant Payment Settlement, центробанковская платформа расчёта мгновенных платежей в~евро под управлением Eurosystem; работает 24/7/365 в~деньгах центрального банка \\
  \textbf{TLV} & Tag-Length-Value, формат кодирования данных EMV (BER-TLV) \\
  \textbf{TMK} & Terminal Master Key, мастер-ключ POS-терминала; под ним хранятся остальные ключи устройства \\
  \textbf{TMS} & Terminal Management System, система управления парком терминалов \\
  \textbf{TPK} & Terminal PIN Key, ключ шифрования PIN внутри POS-терминала \\
  \textbf{TPP} & Third Party Provider, сторонний поставщик услуг в~открытом банкинге PSD2 \\
  \textbf{TPAP} & Third-Party Application Provider, небанковское приложение в~UPI (PhonePe, Google Pay, Paytm), работающее через PSP-bank; для одного TPAP действует потолок 30~\% объёма \\
  \textbf{TPR} & Transaction Processing Rules, документ Mastercard, описывающий правила формирования авторизационных сообщений \\
  \textbf{TPS} & Transactions Per Second, операций в~секунду; целевая метрика производительности \\
  \textbf{TRA} & Transaction Risk Analysis, освобождение от SCA на основе риск-анализа эквайера (PSD2) \\
  \textbf{TRID} & Token Requestor ID, идентификатор запросчика токена в~EMVCo \\
  \textbf{TSP} & Token Service Provider, провайдер токенизации \\
  \textbf{TVR} & Terminal Verification Results, результаты проверок терминала по операции EMV \\
  \textbf{UBO} & Ultimate Beneficial Owner, конечный бенефициарный владелец (KYB) \\
  \textbf{UPI} & Unified Payments Interface, индийский A2A-контур мгновенных платежей \\
  \textbf{UZCARD} & Узбекская национальная платёжная система \\
  \textbf{VAA} & Visa Advanced Authorization, риск-сервис Visa, оценивающий вероятность фрода в~авторизации \\
  \textbf{VAMP} & Visa Acquirer Monitoring Program, программа Visa мониторинга эквайеров по фроду и~chargeback \\
  \textbf{VAU} & Visa Account Updater, сервис автоматического обновления реквизитов карт у~мерчантов \\
  \textbf{VOP} & Verification of Payee, обязательная по IPR проверка соответствия имени получателя и~IBAN перед SCT/SCT Inst-инициацией; схема EPC, действует с~5~октября 2025~года \\
  \textbf{VPA} & Virtual Payment Address, идентификатор получателя в~UPI вида \texttt{name@psp}, отображающий пользователя на банковский счёт без раскрытия реквизитов \\
  \textbf{VCAS} & Visa Consumer Authentication Service, ACS-как-сервис Visa для эмитентов в~3-D~Secure \\
  \textbf{VCPS} & Visa Contactless Payment Specification, спецификация бесконтактных платежей Visa \\
  \textbf{VCR} & Visa Claim Resolution, регламент работы со~спорами Visa \\
  \textbf{VDMP} & Visa Dispute Monitoring Program, программа Visa мониторинга мерчантов по доле споров \\
  \textbf{VEPS} & Visa Easy Payment Service, упрощённый бесконтактный сценарий с~ограничением суммы и~без CVM \\
  \textbf{VFMP} & Visa Fraud Monitoring Program, программа Visa мониторинга мерчантов по доле фрода \\
  \textbf{VROL} & Visa Resolve Online, инструмент Visa для работы со~спорами и~обмена доказательствами \\
  \textbf{VT} & Virtual Terminal, ввод реквизитов оператором ТСП через защищённую страницу (PCI SAQ~C-VT) \\
  \textbf{VTS} & Visa Token Service, токенизационный сервис Visa \\
  \textbf{WAF} & Web Application Firewall, межсетевой экран уровня приложений \\
  \textbf{XS2A} & Access to Account, фреймворк доступа к~счёту в~PSD2/Berlin Group \\
  \textbf{ZMK} & Zone Master Key, мастер-ключ зоны; передаётся между организациями (банк, процессор, сеть) для защищённого обмена другими ключами \\
\\
  \textbf{АСВ} & Агентство по страхованию вкладов \\
  \textbf{АФТ} & Ассоциация ФинТех, разработчик отраслевых стандартов открытых API в~России \\
  \textbf{БИК} & Банковский идентификационный код, девятизначный код участника платёжной системы Банка России \\
  \textbf{БПА} & Банковский платёжный агент (161-ФЗ) \\
  \textbf{БЭСП} & Банковские электронные срочные платежи, исторический сервис ВРСС Банка России (выведен из~эксплуатации в~2018~г.) \\
  \textbf{ГРБС} & Главный распорядитель бюджетных средств; орган, определённый ст.~158 Бюджетного кодекса РФ, имеющий право распределять бюджетные ассигнования и~лимиты обязательств между подведомственными получателями \\
  \textbf{ДБО} & Дистанционное банковское обслуживание (мобильный банк, интернет-банк и~др.) \\
  \textbf{ЕАГ} & Евразийская группа по противодействию легализации преступных доходов и~финансированию терроризма \\
  \textbf{ЕБС} & Единая биометрическая система \\
  \textbf{ЕЦБ} & Европейский центральный банк \\
  \textbf{ЕЭЗ} & Европейская экономическая зона \\
  \textbf{МБВ} & Мир Бюджетные Выплаты, продуктовый стандарт карт Мир для бюджетных получателей \\
  \textbf{МВК} & Межведомственная комиссия по противодействию финансированию терроризма \\
  \textbf{Наличные на кассе} & Сервис ПС~Мир \enquote{покупка с~выдачей наличных} (аналог PWCB); для карт, эмитированных после 10.04.2023, лимит 5\,000~рублей, обязательны онлайн-PIN или CDCVM \\
  \textbf{НКО} & Небанковская кредитная организация \\
  \textbf{НПС} & Национальная платёжная система Российской Федерации (161-ФЗ) \\
  \textbf{НСПК} & Национальная система платёжных карт, оператор карты Мир и~СБП \\
  \textbf{ОД} & Отмывание доходов, полученных преступным путём; вместе с~ФТ образует ПОД/ФТ \\
  \textbf{ОПКЦ} & Операционный и~платёжный клиринговый центр НСПК, обрабатывающий внутрироссийские транзакции международных схем \\
  \textbf{ОСР} & Оператор сервиса рассрочки, статус по~53-ФЗ от~2024~г. \\
  \textbf{ОУД4} & Оценочный уровень доверия~4 по ГОСТ Р~ИСО/МЭК~15408 (аналог Common Criteria EAL4) \\
  \textbf{ОУПИ} & Оператор услуг платёжной инфраструктуры (операционный, платёжный клиринговый, расчётный центры) \\
  \textbf{ОЦ} & Операционный центр в~архитектуре НПС (см.~ОУПИ) \\
  \textbf{ПКС} & Платформа коммерческих согласий Банка России \\
  \textbf{ПКЦ} & Платёжный клиринговый центр в~архитектуре НПС (см.~ОУПИ) \\
  \textbf{ПОД} & Противодействие отмыванию доходов; вместе с~ФТ образует ПОД/ФТ \\
  \textbf{ППС} & Перспективные платёжные сервисы Банка России (преемник БЭСП) \\
  \textbf{ПС} & Платёжная система \\
  \textbf{ПУ} & Поставщик услуг в~российском открытом банкинге (СТО~БР ФАПИ.СЕК); аналог ASPSP \\
  \textbf{ПЭП} & Политически значимое лицо (politically exposed person) \\
  \textbf{РСБУ} & Российские стандарты бухгалтерского учёта \\
  \textbf{РЦ} & Расчётный центр в~архитектуре НПС (см.~ОУПИ) \\
  \textbf{СБП} & Система быстрых платежей Банка России \\
  \textbf{СКЗИ} & Средство криптографической защиты информации, сертифицированное ФСБ России \\
  \textbf{СПУ} & Сторонний поставщик услуг в~российском открытом банкинге; аналог TPP \\
  \textbf{СПФС} & Система передачи финансовых сообщений Банка России \\
  \textbf{ССНП} & Сервисы срочного и~несрочного перевода в~платёжной системе Банка России (Положение Банка России от~25.07.2022 №~802-П) \\
  \textbf{СТО} & Стандарт организации; в~открытом банкинге~-- СТО~БР (Банка России) \\
  \textbf{ТК} & Технический комитет Росстандарта; ТК~26 разрабатывает криптографические стандарты \\
  \textbf{УФЭБС} & Унифицированный формат электронных банковских сообщений Банка России \\
  \textbf{ФАПИ} & Финансовые API; русское обозначение профиля FAPI в~стандартах СТО~БР \\
  \textbf{ФАПИ.СЕК} & Профиль безопасности открытых API Банка России (СТО~БР ФАПИ.СЕК) \\
  \textbf{ФАС} & Федеральная антимонопольная служба \\
  \textbf{ФинЦЕРТ} & Центр мониторинга и~реагирования на компьютерные атаки в~кредитно-финансовой сфере Банка России \\
  \textbf{ФТ} & Финансирование терроризма; вместе с~ОД образует ПОД/ФТ \\
  \textbf{ФЭС} & Формализованное электронное сообщение (115-ФЗ; обмен с~Росфинмониторингом) \\
  \textbf{ЦБ~РФ} & Центральный банк Российской Федерации (Банк России) \\
  \textbf{ЦНФС} & Центр новых финансовых сервисов, дочернее общество Сбербанка \\
  \textbf{ЭДС} & Электронные денежные средства (161-ФЗ) \\
  \textbf{ЭЛКАРТ} & Кыргызская национальная платёжная система; см.~также ELCARD \\
  \textbf{ЭСП} & Электронное средство платежа (161-ФЗ) \\
\end{longtable}
\endgroup
